\section{Introduction}\label{sec:introduction}

\noindent The Electron-Ion Collider will use the ePIC detector to study the structure of nucleons and nuclei through high-luminosity electron-proton and electron-ion collisions. In this environment, forward and backward calorimetry provide important information about neutral particles and event activity outside the central tracking region. The negative-eta hadronic calorimeter (nHCal) is a backward sampling calorimeter intended to improve hadronic coverage in the electron-going region. Optimizing its response to low-energy neutrons is therefore relevant for understanding the acceptance and hit-level performance of the backward hadronic calorimetry region.

\noindent Within that detector context, hadronic calorimeters are typically designed under a limited set of geometric and material constraints, and previous optimization studies often 
focus on choices such as total layer count and absorber-to-scintillator ratio. In these studies, the sampling structure is commonly treated as uniform across depth. This motivates a 
natural question: whether allowing the absorber and scintillator thicknesses to vary non-uniformly can improve particle detection efficiency.

\noindent This note studies that question for low-energy neutron detection at the nHCal. The reference design is a uniform 10-layer 
sampling calorimeter with identical absorber and scintillator thickness in every segment. The alternative explored here is a nonuniform longitudinal parameterization in 
which the calorimeter is divided into three segments, each with its own absorber and scintillator thickness. This provides a coarse variation of the sampling structure while 
avoiding the less realistic case in which every layer is independently optimized.

\noindent Because the resulting parameter space is continuous and multi-dimensional, a brute-force scan is inefficient. The optimization is thus approached with a
surrogate-guided workflow: simulated detector responses are used to train a LightGBM regression model, and that model is then used to guide subsequent proposals of high-performing candidates.

\noindent The purpose of the surrogate-guided workflow is therefore twofold. First, it provides a practical way to search a continuous geometry space without simulating a dense grid of all possible thickness combinations. Second, it tests whether the simulated response contains enough structure for such a workflow to identify useful geometry trends rather than only statistical fluctuations. The main comparison throughout the note is between the uniform nHCal baseline and the best-performing nonuniform geometries identified by this procedure.
